\section{What This Thesis Covers}
\label{sec:intro-scope}

The thesis covers development from the first open-frame laboratory prototype through the integrated V2 board and the reduced V3 edge-card design. The manufactured V1.0 board was tested after bring-up modifications, and later corrections were recorded in schematic V1.1. The V2 design and its preliminary powered tests form the next stage. The hardware chapter identifies the drawing and fitted configuration used at each stage.

The revised plan reports an earlier V2 result near \SI{85}{\watt}. The matched electrical measurements analyzed here were made on the modified V1.0 assembly, reaching \SI{79.89}{\watt}. Later V2 tests examined tank changes and heating. Section~\ref{sec:revised-plan-criteria} sets out the remaining V3 layout checks and the planned matched V2--V3 comparison.

The output section was tested electrically at fixed settings of \SI{5}{\volt}, \SI{9}{\volt}, \SI{15}{\volt}, and \SI{20}{\volt}. PPS characterization remains part of the revised plan. USB-IF protocol and interoperability tests, product-safety qualification, surge and ESD qualification, conducted and radiated emissions tests, and long-duration reliability tests remain outside the completed test set.

The planar components were assessed in the installed converter: their connections, mechanical support, spacing, thermal environment, and effect on operation. Core selection, winding optimization, the original FEM models, and electromagnetic design belong to the supplied work.

The supplied analysis and LTspice circuits allow circuit relationships and approximate component values to be checked. Chapter~\ref{ch:simulation} describes the calculations and the changes needed to model the final board and firmware.

Long-duration thermal testing is still needed to determine the continuous output rating.
